\chapter{热插拔与驱动自动装载}
\label{chap:hotplug}

第~\ref{chap:drivers} 章的 USB 键盘、MT7601U 都是“开机时插好、启动时枚举一次”的设备。本章讨论设备在运行中出现和消失：Hub 用自己的 interrupt-IN 端点报告端口变化，主机据此对新插入的设备重新枚举、拔出时拆除整棵子树；当总线上找不到匹配的驱动时，设备模型生成 MODALIAS 字符串，经 \code{request\_module()} 从根文件系统的 \file{/lib/modules} 装载对应的 \file{.ko}。前者依赖第~\ref{chap:drivers} 章的 DWC2 主机通道、split 事务和 URB（\ref{sec:usb-hotplug} 节），后者依赖设备模型的 match/probe 与 \code{modalias} 钩子（\ref{sec:devmodel} 节），最后用一份 Logitech 复合接收器的真机日志把两者串起来（\ref{sec:usb-mouse-autoload} 节）。

\section{Hub 中断驱动的 USB 热插拔}
\label{sec:usb-hotplug}

键盘按键报告和 USB 拓扑变化是两种不同的事件。按键数据由键盘自己的
interrupt-IN URB 接收；热插拔则由 LAN951x Hub 的 interrupt-IN endpoint
通知主机。当前 DWC2 为 Hub 的长期 URB 固定使用 channel 7，为唯一 HID
数据 URB 使用 channel 6。Hub 空闲时会对 interrupt-IN 返回 NAK，HCD 按
\code{bInterval} 重新调度请求；只有端口变化位图非零时才唤醒 Hub 事件
worker，不再每隔固定时间轮询所有端口的 EP0。

整个“插入键盘→枚举→可按键”的过程，可以串成一条从端口变化到第一个按键的总线时间线（图~\ref{fig:usb-hid-enum-sequence}）：开机时四个端口都空载，控制器把 ch7 挂在 Hub 的 interrupt-IN 上等待热插拔；插入 port 3 之后由端口变化中断触发 reprobe，低速键盘的 EP0 枚举全部经 Hub 的 TT 拆成 SSPLIT/CSPLIT；驱动绑定完成后，周期 interrupt-IN 固定交给 ch6，直到某个按键的 HID 报告变成用户空间的事件。

\begin{figure}[p]
  \centering
  \resizebox{\linewidth}{!}{%
  % 图（fig:usb-hid-enum-sequence）：DWC2 经外部 Hub 枚举低速 HID 键盘的总线时序（热插拔版）。
% 五个阶段对应真机串口日志：空载扫描 → ch7 端口变化中断 → SSPLIT/CSPLIT 枚举
% → ch6 武装 → 第一个按键。阶段配色沿用 mermaid 原稿的底色。
\definecolor{seqstgA}{RGB}{240,248,255}
\definecolor{seqstgB}{RGB}{255,245,238}
\definecolor{seqstgC}{RGB}{245,255,245}
\definecolor{seqstgD}{RGB}{255,240,245}
\definecolor{seqstgE}{RGB}{255,250,200}
\begin{tikzpicture}[x=1mm,y=-1mm,font=\scriptsize,
    participant/.style={draw=structurecolor!75!black,fill=structurecolor!12,
      rounded corners=1.5pt,align=center,text width=24mm,minimum height=9mm},
    lifeline/.style={gray!65,dashed,line width=.45pt},
    seqcall/.style={-{Stealth[length=1.6mm]},line width=.55pt,
      draw=structurecolor!85!black},
    seqret/.style={-{Stealth[length=1.6mm]},dashed,line width=.5pt,
      draw=structurecolor!75!black},
    msg/.style={fill=white,inner sep=1pt,align=center,font=\tiny},
    note/.style={draw=gray!55,fill=yellow!16,rounded corners=1.5pt,inner sep=1.5pt,
      align=center,font=\tiny,text=gray!70!black},
    stage/.style={anchor=west,align=left,font=\scriptsize\bfseries,
      text=structurecolor!80!black},
    loopbox/.style={draw=gray!60,dashed,line width=.4pt,rounded corners=2pt},
    looplab/.style={draw=gray!60,fill=white,rounded corners=1.5pt,inner sep=1.5pt,
      font=\tiny,text=gray!70!black,anchor=north west}]
  \node[participant] at (0,0) {DWC2 控制器\\Host／调度器};
  \node[participant] at (46,0) {外部集线器\\addr 2，4 ports};
  \node[participant] at (92,0) {USB 键盘\\低速设备};
  \node[participant] at (138,0) {内核 \code{usbkbd}\\Input 子系统};
  \foreach \x in {0,46,92,138}
    \draw[lifeline] (\x,5) -- (\x,213);
  \begin{scope}[on background layer]
    \fill[seqstgA,rounded corners=2pt] (-17,8) rectangle (156,48);
    \fill[seqstgB,rounded corners=2pt] (-17,50) rectangle (156,76);
    \fill[seqstgC,rounded corners=2pt] (-17,78) rectangle (156,134);
    \fill[seqstgD,rounded corners=2pt] (-17,136) rectangle (156,160);
    \fill[seqstgE,rounded corners=2pt] (-17,162) rectangle (156,212);
  \end{scope}

  % 阶段 1：冷启动，总线上没有可枚举的子设备
  \node[stage] at (-15,10) {阶段 1：系统冷启动与空载扫描};
  \draw[seqcall] (0,17) -- (0,20.5);
  \node[msg,anchor=west] at (2,18.7) {HCD Power On，开启 channel IRQ};
  \draw[seqcall] (0,24) -- node[msg,above]{枚举 Root Hub 及外部集线器（addr 2）} (46,24);
  \draw[seqcall] (0,29) -- node[msg,above]{扫描 port 1／2／3／4 状态} (46,29);
  \draw[seqret] (46,36) -- node[msg,above=1.5pt]{四个端口状态均为 \code{0x100}（空载）\\
    日志 \code{no supported USB child}，转入待机} (0,36);
  \draw[seqcall] (0,43) -- (0,46.5);
  \node[msg,anchor=west] at (2,44.7) {分配 ch7：开启事件驱动热插拔监听（hotplug v2）};

  % 阶段 2：Hub 的 interrupt-IN 报告端口变化
  \node[stage] at (-15,52) {阶段 2：热插拔事件触发（用户把键盘插入 port 3）};
  \draw[seqret] (46,59) -- node[msg,above=1.5pt]{\textbf{[ch7 IRQ]} 端口状态改变
    \code{change=8}，\code{connected=2$\rightarrow$a}\\收到中断，触发 reprobe 重新探测} (0,59);
  \draw[seqcall] (0,67) -- node[msg,above]{\code{RESET\_PORT(3)}：复位端口} (46,67);
  \draw[seqret] (46,72) -- node[msg,above=1.5pt]{复位完成，\code{portstatus=0x303}\\
    识别为 low speed（1.5\,Mb/s）} (0,72);

  % 阶段 3：低速设备的 EP0 枚举全部走 split transaction
  \node[stage] at (-15,80) {阶段 3：分割事务枚举（split transaction）；\\
    \hspace*{2em}低速设备必须经 Hub 的 TT 执行 SSPLIT／CSPLIT};
  \draw[seqcall] (0,89) -- node[msg,above]{SSPLIT：\code{GET\_DESCRIPTOR(Device)}} (46,89);
  \draw[seqcall] (46,94) -- node[msg,above]{低速 EP0 读设备描述符} (92,94);
  \draw[seqret] (46,99) -- node[msg,above]{CSPLIT：返回 \code{VID=2a7a}，\code{PID=8a47}，
    \code{MPS=8}} (0,99);
  \draw[seqcall] (0,105) -- node[msg,above]{SSPLIT：\code{SET\_ADDRESS(4)}} (46,105);
  \node[note,anchor=west] at (94,110) {设备地址更新为 4};
  \draw[seqcall] (0,116) -- node[msg,above]{SSPLIT：\code{GET\_DESCRIPTOR(HID／Config)}} (46,116);
  \draw[seqret] (46,121) -- node[msg,above]{CSPLIT：返回 HID 报告描述符} (0,121);
  \draw[seqcall] (0,126) -- node[msg,above]{SSPLIT：\code{SET\_CONFIGURATION}（激活设备）} (46,126);
  \draw[seqcall] (46,131) -- node[msg,above]{激活 HID 接口} (92,131);

  % 阶段 4：类驱动绑定，周期 interrupt-IN 交给 ch6
  \node[stage] at (-15,138) {阶段 4：驱动绑定与轮询武装};
  \draw[seqcall] (138,144) -- (138,147.5);
  \node[msg,anchor=east] at (136,145.7) {识别为 Boot Keyboard，注册 input 设备 \code{event0}};
  \draw[seqcall] (0,151) -- (0,154.5);
  \node[msg,anchor=west] at (2,152.7) {分配 ch6：\code{addr=4}，\code{EP=1}，
    \code{interval=10\,ms}；\\ch6 武装完毕（armed），ch7 继续在后台监听 Hub 热插拔};
  \draw[seqcall] (138,158) -- node[msg,above]{提交 URB，开始监听按键} (0,158);

  % 阶段 5：轮询命中一次真实的按键
  \node[stage] at (-15,164) {阶段 5：真实按键捕获};
  \draw[loopbox] (-14,168) rectangle (110,183);
  \node[looplab] at (-13.4,168.4) {loop／每 10\,ms 轮询一次};
  \draw[seqcall] (0,174) -- node[msg,above]{\textbf{[ch6]} SSPLIT IN：有按键吗？} (46,174);
  \draw[seqcall] (46,180) -- node[msg,above]{低速 IN 令牌} (92,180);
  \draw[seqret] (92,188) -- node[msg,above=1.5pt]{用户按下某个键（例如 \code{f}），返回 8 字节 HID 报告\\
    \code{0x00, 0x00, 0x0f, 0x00, 0x00, 0x00, 0x00, 0x00}} (46,188);
  \draw[seqret] (46,195) -- node[msg,above]{\textbf{[ch6]} CSPLIT：传输完成} (0,195);
  \draw[seqcall] (0,200) -- node[msg,above]{\textbf{[ch6 IRQ]} 上报
    \code{data=0,0,f,0,0,0,0,0}} (138,200);
  \draw[seqcall] (138,206) -- (138,209.5);
  \node[msg,anchor=east] at (136,207.7) {\code{usbkbd} 解析：\code{mod=0}，
    \code{keys=f}（HID Usage \code{0x0f}）；\\Input 子系统生成按键事件，发送至用户空间};
\end{tikzpicture}

  }
  \caption{DWC2 经外部 Hub 热插拔枚举低速 HID 键盘、直到收到第一个按键的总线时序。
  ch7 只报告 Hub 的端口变化，worker 收到中断才 reprobe；EP0 枚举用 SSPLIT/CSPLIT
  经 Hub 的 TT 转发；驱动绑定后 ch6 以 10\,ms 周期轮询按键，HID 报告最终由输入子
  系统发布为事件}
  \label{fig:usb-hid-enum-sequence}
\end{figure}

\begin{txtcode}
Hub interrupt endpoint
  -> DWC2 channel 7 DMA receives change bitmap
  -> IRQ 9 records completion and queues hub_event_work
  -> dwc2_hub_wq examines changed ports only
  -> GET_STATUS / CLEAR_FEATURE for those ports
  -> compare connected bitmap
\end{txtcode}

硬中断只处理 host-channel 完成并排入工作项，不在 IRQ 中进行 Hub control
传输。独立的 \code{dwc2\_hub\_wq} 根据 bitmap 中标记的端口发送
\code{GET\_STATUS}，读取连接状态和变化位，再用 \code{CLEAR\_FEATURE}
清除已处理的 port change feature。比较新旧 connected bitmap 后，只有拓扑
确实变化才启动重枚举；因此空闲 NAK 不会造成无意义的全端口扫描。

设备拔出时，类驱动和 input/evdev 的资源按逆序回收：

\begin{txtcode}
usb_device_unregister(old_child)
  -> usbhid remove()
  -> mark device disconnected
  -> usb_kill_urb()
       -> stop/mask the HID host channel
       -> wait for completion callback
  -> cancel_work_sync()
  -> input_unregister_device()
       -> remove event-minor registration
       -> wake blocked event readers
  -> release driver-owned references
\end{txtcode}

注销 event 设备会唤醒阻塞在 \code{read()} 的客户端，剩下的语义与
\ref{sec:usb-input} 小节一致：旧 fd 先读完已提交的完整帧，内存等最后一个
引用释放。热插拔在这里第一次真正考验了这套引用计数——拔掉的瞬间
\code{evtest} 往往还活着，它必须看到断开错误而不是继续读到一个已经不属于
自己的设备。

\paragraph{channel 7 怎么保证自己活着.} Hub 的 interrupt-IN 是有周期要求的：
\code{dwc2\_hub\_interrupt\_start()} 提交 ch7 时按当前 \reg{HFNUM} 的奇偶选择
\code{HCCHAR.ODDFRM}，把事务安排在下一个 USB frame；同时在
\code{dwc2\_hub\_wq} 上用 \code{mod\_delayed\_work()} 登记一个 \code{bInterval}
期限的看门狗。正常的 NAK 完成 IRQ 会把这条期限往后推；如果控制器复位后的
第一笔事务根本没有产生完成 IRQ，期限到期就会重新启动 ch7。它只恢复 Hub 的
interrupt URB，不做任何 \code{GET\_STATUS}，因此不是退回旧的端口轮询。
新一代监控还会先等一个 \code{bInterval} 再首次提交，确保触发重枚举的那只旧
URB 的 completion 已经完全退出。下面这三行里后两行各只打印一次，合起来正好
证明 delayed work 和 DWC2 host channel 两段链路都在运行：

\begin{txtcode}
dwc2: event-driven hub hotplug v2 scheduled on channel 7
dwc2: hub ch7 scheduler alive active=0 urb-state=0
dwc2: hub ch7 first IRQ hcint=... hctsiz=...
\end{txtcode}

一次真机热插 port 3 键盘时，连接位图从 \code{0x2} 变为 \code{0xa}：
\code{0x2} 表示 port 1 已连接，\code{0xa} 则表示 port 1 和 port 3 已连接，
位图的 bit 号就是 root hub 的端口号，所以新增设备来自 port 3。Hub
interrupt URB 通知 worker 后，可看到类似日志：

\begin{txtcode}
dwc2: hub interrupt-IN ep=1 mps=1 interval=12 ports=4
dwc2: hub IRQ change=8 connected=2->a; reprobe
dwc2: hub port 3 status=301 change=1
dwc2: hub port 3 after-reset status=303 speed=low
dwc2: hub port=3 child class=0 vid=2a7a pid=8a47 mps=8
input: USB HID Boot Keyboard registered as event0
usbkbd: IRQ-driven boot keyboard ready ep=1 mps=8 interval=10
\end{txtcode}

\code{status=0x301} 表示设备已连接、端口已供电且设备为低速；
reset 后的 \code{0x303} 还包含端口使能状态。若插入的是 Boot Mouse，USB
core 会绑定鼠标驱动并注册相应 input 设备。枚举和 probe 日志只证明新设备
已绑定、接收 URB 已提交；要证明实际按键路径恢复，还应看到 channel 6 的
completion 和首份有效 HID report。

\paragraph{拔出后的短暂 \code{ch6 error} 是正常窗口.} 键盘刚拔掉、Hub 的状态
端点还没把变化报上来之前，常见这几行：

\begin{txtcode}
dwc2: keyboard ch6 error intr=82 hctsiz=... hcsplt=...
usbkbd: channel 6 transfer error; rearming
\end{txtcode}

\code{0x82} 是 \code{CHHLTD | XACTERR}：ch6 仍在访问一个已经离线的低速端点，
这在“Hub 状态端点的下一个周期到来之前”是必然窗口，不是键盘驱动的残留问题。
判断依据是后续日志：只要接着出现 \code{hub IRQ change} 和
\code{... event0 disconnected}，remove 生命周期就已经收敛。反过来，如果只有
\code{ch6 error} 而没有 Hub 的变化事件，那要查的是 ch7 那条链路。

\paragraph{DMA 缓冲的所有权.} 插回键盘时如果连续打印

\begin{txtcode}
dwc2: control data failed req=0 addr=1 result=-1
\end{txtcode}

其中 \code{req=0, addr=1} 是事件 worker 向根 Hub 发端口
\code{GET\_STATUS}，并不是周期拓扑轮询出错了。这里的根因很隐蔽：早期实现把
4 字节的端口状态数组放在 worker 的内核栈上，又直接把它交给 DWC2 做 DMA。
缓存维护按 64 字节 cache line 向外取整，\emph{invalidate} 时会把同一条
cache line 里的其它局部变量一起覆盖掉；而且这样一个短生命周期的栈对象本来
就不该成为 DMA 目标。现在端口状态用独立的、静态存续且 64 字节对齐的
\code{dwc2\_hub\_port\_status[64]}，控制传输仍只接收前 4 字节；Hub 状态位图、
端口状态缓冲和枚举用的 \code{ctrl\_buf} 三者彼此分离，DMA 与缓存的归属才不会
互相污染。这也是本章反复出现的那条规则的又一次兑现：交给 DMA 的缓冲区必须
由传输的 owner 独占，并且活得比传输久。

变化之后的恢复过程是这样的：HCD 先注销那个唯一的发布用
\code{usb\_child}，停止并屏蔽相关 channel、清理待处理中断，
然后重新初始化 DWC2 并扫描 Hub。新设备按地址 0 开始枚举，读取 Device
Descriptor 前 8 字节确定 EP0 最大包长，设置 USB 地址、读取配置描述符并
选择 Boot Keyboard 或 Boot Mouse 接口；随后 USB bus 匹配驱动，执行 probe，
重新注册 input 设备并提交对应 HID URB。对低速键盘，地址路由还需记录父 Hub、
端口和速度，以便 EP0 与 channel 6 的传输正确使用 split transaction。

\paragraph{与 Linux 还差什么.} 把这套驱动和真正的 Linux USB 栈对比，剩下的差距是明确的几条：

\begin{itemize}
  \item \textbf{重枚举的粒度。}Linux 只增删发生变化的 port 和 device；这里的一次拓扑变化会注销旧设备、停止并屏蔽相关 channel，然后重新初始化整个 DWC2。嵌套 Hub 的下游端口变化甚至不会出现在根 Hub 的连接位图里。
  \item \textbf{host channel 的分配。}ch7 固定给 Hub 的 interrupt URB；ch6 不再独占某一只 HID 数据 URB，而是由 \code{interrupt\_rx[DWC2\_INT\_REQUESTS]} 数组在多个 interrupt-IN 请求之间轮转仲裁（同一时刻只有一个请求拥有 ch6）。这让键盘与鼠标可以共处一只 channel，但它是时分复用，不是 Linux 那种按设备/端点动态分配多 channel 的通用调度。
  \item \textbf{周期调度表。}空闲 NAK 已经用 \code{retry\_work} 加 \code{dwc2\_periodic\_wq} 按 \code{bInterval} 重排，不再占着公共 worker 自旋；但 10\,ms 的 delayed-work 粒度只适合这只低速键盘，通用的 HCD 需要以 USB frame/SOF 为不同端点各自维护周期表。
  \item \textbf{硬中断里的忙等。}为了满足 split 的 125\,$\mu$s microframe 期限，快速路径最多在 IRQ 里等 500\,$\mu$s 并重启 channel。Linux DWC2 用更完整的 SOF/periodic 状态机把这部分忙等压得更低。
  \item \textbf{多设备发布。}Boot HID 接口现在按接口逐个发布独立的 \code{usb\_device}（\code{usb\_hid\_child[]}，最多 \code{DWC2\_HID\_CHILDREN}=4 个），各自保存地址、接口三元组与拔出生命周期，因此复合接收器的键盘接口与鼠标接口能被不同驱动分别认领。非 HID 的 CDC/bulk 设备（如 MT7601U、usbnet）仍共用单个 \code{usb\_child}，与 HID 子设备尚未统一成一张按 Hub 端口排列的数组。
\end{itemize}

\section{复合 HID 接口与鼠标驱动自动装载（真机日志 046d:c534）}
\label{sec:usb-mouse-autoload}

\ref{sec:usb-debug} 节的那台接收器是 low-speed、单键盘接口。真机上更常见的是 Logitech 无线接收器
\code{046d:c534} 这类 \textbf{full-speed 复合 HID}：同一个 USB 地址上接口 0 是 Boot
Keyboard（\code{3/1/1}）、接口 1 是 Boot Mouse（\code{3/1/2}）。这一节沿一份真实串口日志
把“插入—split 枚举—双接口分别登记—键盘即时绑定—鼠标经 MODALIAS 从 ext2 根自动装载
\texttt{.ko}—拔出触发整树重枚举”走一遍，并说明 \file{dwc2.c} 的分接口发布
与 \file{module.c} 的自动装载链在这里如何被验证。

\paragraph{拓扑与速度.} 连接位图 \code{connected=2->a}（bit1\,|bit3）说明新增设备来自
port 3；复位后 \code{status=103} 的速度域是 \textbf{full-speed}。这是与低速键盘最容易混淆
的一点：full-speed 设备在高速 Hub 之后\emph{同样}要经 TT 做 split 事务，但
\reg{HCCHAR} 的 \code{LSPDDEV} 位\textbf{不置}（只有 low-speed 才置）。表~\ref{tab:usbmouse-regs}
用本日志的真实寄存器值把这条规则摊开。

\begin{table}[htbp]
  \centering
  \caption{鼠标接收器（046d:c534）日志中的报文解码}
  \label{tab:usbmouse-regs}
  \small
  \begin{tabularx}{\linewidth}{@{}lX@{}}
    \toprule
    样本 & 含义 \\
    \midrule
    \reg{HCSPLT}$=$\code{0x8000c083} & \code{SPLTENA}、XACTPOS$=$ALL、Hub 地址 1、端口 3；\code{COMPSPLT}$=0$ 即当前是 SSPLIT（错误行里的 \code{0x8001c083} 则 \code{COMPSPLT}$=1$，为 CSPLIT） \\
    \reg{HCCHAR}$=$\code{0x810c8808} & 地址 4、端点 1、IN、interrupt 类型、最大包长 8，\code{LSPDDEV}$=0$：full-speed \\
    \reg{HCINT}$=$\code{0x22} & \code{ACK | CHHLTD}：TT 接受 SSPLIT \\
    \reg{HCINT}$=$\code{0x23} & \code{XFERCOMPL | ACK | CHHLTD}：CSPLIT 取回完整数据 \\
    \reg{HCINT}$=$\code{0x0a} & \code{STALL | CHHLTD}：\code{control status failed req=10} 的 STALL（\code{req=10}$=$\code{0x0a}$=$SET\_IDLE） \\
    \reg{HCINT}$=$\code{0x82} & \code{XACTERR | CHHLTD}：拔出前 \code{HID ch6 error} \\
    \bottomrule
  \end{tabularx}
\end{table}

一次 EP0 完成事务的 SSPLIT/CSPLIT 时序，与图~\ref{fig:usb-split} 形状相同（只是本例是
full-speed），日志里的 \code{hfnum} 记录了两拍的帧号；interrupt-IN 阶段的
\code{dwc2 ch6 final p=0/1 ...} 打印则显示 ch6 在键盘与鼠标两只 URB 之间轮转：\code{p} 是
\code{req->complete\_split}（0=SSPLIT、1=CSPLIT），\code{issue}/\code{ssack}/\code{age} 分别是发出帧、
TT 接受 SSPLIT 帧与距发出的帧数，是 CSPLIT 有效窗口的判据。

\paragraph{分接口发布.} \file{dwc2.c} 的枚举在 \code{SET\_CONFIGURATION} 后遍历整段配置描述符，
对每个带 interrupt-IN 端点的 Boot HID 接口各建一个 \code{usb\_device}（\code{usb\_hid\_child[]}），
逐个 \code{usb\_device\_register()}。日志因此出现两行：

\begin{txtcode}
dwc2: HID boot keyboard intf 0 handed to usb bus
dwc2: HID boot mouse   intf 1 handed to usb bus
\end{txtcode}

接口 0 的 \code{3/1/1} 命中内置 \file{usbkbd}，同步 probe 并占用 ch6；接口 1 的
\code{3/1/2} 没有任何已注册驱动匹配（内置键盘只认 \code{protocol=1}），于是设备 core 走
未匹配分支：用 bus 的 \code{modalias()} 生成别名并 \code{request\_module()}。

\begin{table}[htbp]
  \centering
  \caption{Boot Mouse 的 MODALIAS 与 HID control、报告报文}
  \label{tab:usbmouse-packets}
  \small
  \begin{tabularx}{\linewidth}{@{}lX@{}}
    \toprule
    项 & 内容 \\
    \midrule
    MODALIAS & \code{usb:c03s01p02}（class/subclass/protocol 的十六进制）；静态 alias 表映射到 \code{usbmouse.ko} \\
    SET\_PROTOCOL & SETUP：\code{0x21 0x0b}，\code{wValue=0, wIndex=interface, wLength=0}；失败才放弃 probe \\
    SET\_IDLE & SETUP：\code{0x21 0x0a}，同上；失败不致命 \\
    Boot Mouse 报告 & 4 字节：byte0 按钮位图（左/右/中），byte1 X、byte2 Y、byte3 滚轮（均为 signed char）；\code{actual\_length>=3} 即接受 \\
    输入事件 & 按钮变化 \code{input\_report\_key(BTN\_*)}；位移 \code{input\_report\_rel(REL\_X/Y/WHEEL)}；末尾 \code{input\_sync()} \\
    \bottomrule
  \end{tabularx}
\end{table}

\paragraph{自动装载链.} 从设备注册到鼠标驱动绑定，跨越 USB core、设备 core 与模块
加载器：

\begin{txtcode}
usb_device_register(接口1)  ->  device_register()  ->  无驱动匹配
  ->  usb_modalias() = "usb:c03s01p02"
  ->  request_module(alias)                     (见 device.c)
       ->  命中静态 alias 表 ->  queue_work(kmod_wq, slot)
            ->  kmod_run() ->  kmod_load("/lib/modules/usbmouse.ko")
                 ->  begin_op(); namei(); ilock()
                 ->  ip->iop->read(ip, 0, ko_image, 0, sz)   % ext2 分派，非经典 readi()
                 ->  load_module_image()：模块区分配、符号解析、RELA 重定位
                 ->  run init: usbmouse_driver_init()
                      ->  init_workqueue(usbmouse_wq)        % 日志 "usbmouse_wq worker=0 pid=17"
                      ->  usb_register_driver() -> driver_register()
                           ->  重扫已登记设备 -> usb_match(接口1) 命中
                           ->  usbmouse_probe() -> usb_submit_urb()
                                ->  dwc2_submit_urb() -> dwc2_interrupt_schedule_locked()
                                     若 ch6 空闲则 rx_start，否则排入 interrupt_rx[] 等轮转
\end{txtcode}

两个要点。其一，\code{driver\_register()} 会重新遍历总线上\textbf{已存在}的设备，所以“设备
先登记并异步排队装载、驱动后注册”的次序仍能绑定接口 1，无需手动 \code{insmod}。其二，
\code{kmod\_load()} 里读 \texttt{.ko} 必须走 \code{ip->iop->read}（VFS 分派），不能用经典
xv6fs 的 \code{readi()}：\code{readi()} 的 \code{bmap()->getsb()} 假设设备带 xv6 超级块，在
ext2 根上会直接 \code{panic: getsb: device has no xv6 filesystem}。本日志的验收点正是这行：

\begin{txtcode}
request_module: loading /lib/modules/usbmouse.ko
module: loaded usbhid-mouse (1808 bytes)
\end{txtcode}

出现这两行、且不再有 \code{getsb} panic，说明方案 A 的分接口发布、MODALIAS 自动装载
与 ext2 上的模块读取在真机跑通。

\paragraph{本日志尚未证明的部分.} 按“只承认证据”的原则，这次链路证明了\textbf{自动装载}，
却没有证明\textbf{鼠标数据端到端接收}：全程未出现 \code{usbmouse: IRQ-driven boot mouse
ready}，也没有 \code{usbmouse: first report}；拔出前是连续四次的
\code{dwc2: HID ch6 error intr=82}（\code{XACTERR | CHHLTD}）与 \code{usbkbd: channel 6 transfer
error; rearming}。也就是说 full-speed 复合设备在共享 ch6、经 TT 的 split interrupt-IN
调度仍未收敛，键盘与鼠标争用同一 channel 的时序是主要嫌疑；拔出仍触发
\code{connected=a->2} 的\textbf{整棵拓扑重枚举}，而非只对 port 3 局部增删。下一步诊断应在
\code{usbmouse\_probe()} 各返回前打印失败的 control 请求，并把 ch6 错误路径的
\reg{HCSPLT}/\reg{HCTSIZ}/\reg{HCCHAR} 一并 dump，以区分是 SET\_PROTOCOL 被 STALL、还是
interrupt-IN 的 CSPLIT 窗口错过。
